%% ============================================================================
\section{Reproducibility}
\label{app:repro}

All results derive from a frozen snapshot cataloged in Table~\ref{tab:data}, whose source
URLs come from the data manifest, which also records a SHA-256 digest for every file it
lists. Every stochastic routine is seeded
($\mathrm{SEED}=20230311$).

The hourly coverage is not uniform across coins. USDC is missing 4 of 216 expected hourly
points (17:00~UTC on 13~March, 14:00~UTC on 14~March, 17:00~UTC on 14~March, and 12:00~UTC
on 15~March), DAI 2 of 216 (17:00~UTC on 13~March and 16:00~UTC on 14~March), and USDT 3 of
216 (20:00~UTC on 11~March, 17:00~UTC on 13~March, and 17:00~UTC on 14~March). One hour,
17:00~UTC on 13~March, is missing from all three.

\begin{table}[t]
\centering
\caption{Frozen data snapshot. Every row's count is read directly from the file itself.
The 212/213/214-point spread across USDC/DAI/USDT in row 1 reflects independent gaps in the
upstream aggregator's per-coin coverage, enumerated in the text above.}
\label{tab:data}
\small
\checkwidth{\small% Source: data/MANIFEST.json (files{} + sources[] entries) -- row counts read
%   directly from each listed file, not from MANIFEST prose or estimated:
%   data/price_hourly.csv, data/binance_usdcusdt_1h.csv, data/usdc_supply_daily.csv,
%   data/redemptions_by_wallet.csv, data/usdt_usd_hourly.csv, data/terra_price_hourly.csv,
%   data/n1/kraken_usdcusd/hourly.csv, data/n1/kraken_usdtusd/hourly.csv,
%   data/n1/kraken_usdtusd_2019/hourly.csv, data/n1/onchain/hourly_block_anchors.csv,
%   data/n1/tether2019/disclosure_passages_nyag.json (passages[]),
%   data/n1/paxos2023/disclosure_passages_paxos.json (list),
%   data/n1/paxos2023/usdp_supply_daily.csv,
%   data/n1/labels/dawsbot_eth-labels/accounts.csv,
%   data/n1/labels/brianleect_etherscan-labels/combinedAccountLabels.json (dict keys);
%   results/number4_dai_contagion.json (run_date_utc, DAI backing row's Coverage).
% Width pass 2: see write_tab_sources docstring. Coverage is date-only, read from
%   each file's own timestamp column via _cov_range/_d_only_iso, except Redemptions
%   (constants.WINDOW_START_S/WINDOW_END_S, matching redemptions.py's block_by_time
%   calls) and the two GitHub label rows (no frozen retrieval date exists, so "--"
%   rather than a fabricated one). Series shortened, Source cut to one or two words.
\begin{tabular}{@{}p{2.9cm}p{2.9cm}p{1.5cm}p{2.7cm}@{}}
\toprule
Series & Source & Rows & Coverage \\
\midrule
USDC/USDT/DAI price & DeFiLlama & 639 & 7--16 Mar 2023 \\
USDC/USDT cross & Binance & 131 & 11--17 Mar 2023 \\
USDC supply & DeFiLlama & 12 & 7--18 Mar 2023 \\
Redemptions & Etherscan & 776 & 8--17 Mar 2023 \\
DAI backing & Etherscan & 1 & 11 Mar 2023 \\
USDT/USD price+vol & Coinbase & 217 & 8--17 Mar 2023 \\
Terra price & DeFiLlama & 384 & 7--15 May 2022 \\
\midrule
Kraken USDC/USD & Kraken & 240 & 8--17 Mar 2023 \\
Kraken USDT/USD & Kraken & 240 & 8--17 Mar 2023 \\
Kraken USDT/USD 2019 & Kraken & 624 & 20 Apr--15 May \\
NYAG press release & NYAG & 4 & 25 Apr 2019 \\
Paxos/USDP statement & Paxos & 3+12+1 & 8--14 Mar 2023 \\
Labels (dawsbot) & GitHub & 144{,}378 & -- \\
Labels (brianleect) & GitHub & 23{,}288 & -- \\
Block anchors & Ethereum RPC & 240 & 8--17 Mar 2023 \\
\bottomrule
\end{tabular}
}
% Source: data/MANIFEST.json (files{} + sources[] entries) -- row counts read
%   directly from each listed file, not from MANIFEST prose or estimated:
%   data/price_hourly.csv, data/binance_usdcusdt_1h.csv, data/usdc_supply_daily.csv,
%   data/redemptions_by_wallet.csv, data/usdt_usd_hourly.csv, data/terra_price_hourly.csv,
%   data/n1/kraken_usdcusd/hourly.csv, data/n1/kraken_usdtusd/hourly.csv,
%   data/n1/kraken_usdtusd_2019/hourly.csv, data/n1/onchain/hourly_block_anchors.csv,
%   data/n1/tether2019/disclosure_passages_nyag.json (passages[]),
%   data/n1/paxos2023/disclosure_passages_paxos.json (list),
%   data/n1/paxos2023/usdp_supply_daily.csv,
%   data/n1/labels/dawsbot_eth-labels/accounts.csv,
%   data/n1/labels/brianleect_etherscan-labels/combinedAccountLabels.json (dict keys);
%   results/number4_dai_contagion.json (run_date_utc, DAI backing row's Coverage).
% Width pass 2: see write_tab_sources docstring. Coverage is date-only, read from
%   each file's own timestamp column via _cov_range/_d_only_iso, except Redemptions
%   (constants.WINDOW_START_S/WINDOW_END_S, matching redemptions.py's block_by_time
%   calls) and the two GitHub label rows (no frozen retrieval date exists, so "--"
%   rather than a fabricated one). Series shortened, Source cut to one or two words.
\begin{tabular}{@{}p{2.9cm}p{2.9cm}p{1.5cm}p{2.7cm}@{}}
\toprule
Series & Source & Rows & Coverage \\
\midrule
USDC/USDT/DAI price & DeFiLlama & 639 & 7--16 Mar 2023 \\
USDC/USDT cross & Binance & 131 & 11--17 Mar 2023 \\
USDC supply & DeFiLlama & 12 & 7--18 Mar 2023 \\
Redemptions & Etherscan & 776 & 8--17 Mar 2023 \\
DAI backing & Etherscan & 1 & 11 Mar 2023 \\
USDT/USD price+vol & Coinbase & 217 & 8--17 Mar 2023 \\
Terra price & DeFiLlama & 384 & 7--15 May 2022 \\
\midrule
Kraken USDC/USD & Kraken & 240 & 8--17 Mar 2023 \\
Kraken USDT/USD & Kraken & 240 & 8--17 Mar 2023 \\
Kraken USDT/USD 2019 & Kraken & 624 & 20 Apr--15 May \\
NYAG press release & NYAG & 4 & 25 Apr 2019 \\
Paxos/USDP statement & Paxos & 3+12+1 & 8--14 Mar 2023 \\
Labels (dawsbot) & GitHub & 144{,}378 & -- \\
Labels (brianleect) & GitHub & 23{,}288 & -- \\
Block anchors & Ethereum RPC & 240 & 8--17 Mar 2023 \\
\bottomrule
\end{tabular}

\end{table}

Every reported number and figure regenerates from this snapshot via the following command
sequence: \texttt{python data\_fetch.py verify} $\to$ \texttt{python data\_fetch\_n1.py
verify} $\to$ \texttt{python redemptions.py --from-frozen} $\to$ \texttt{python -m pytest}
$\to$ \texttt{python report.py}. None of these five steps makes a network call or reads
\texttt{ETHERSCAN\_API\_KEY}, and none can overwrite a file already recorded in the data
manifest. \texttt{make fc27} also pins \texttt{SOURCE\_DATE\_EPOCH} to
a fixed constant and the PDF trailer \texttt{/ID} via \texttt{\textbackslash
pdftrailerid\{\}}, so the typeset PDF is byte-reproducible.
Tables~\ref{tab:claimmap} and~\ref{tab:claimmapb} name the file and key behind each.

Acquisition is a separate step and requires network access. Six scripts wrote the frozen
snapshot. \texttt{data\_fetch.py}, \texttt{data\_fetch\_n1.py}, \texttt{second\_event.py},
and \texttt{usdt\_premium.py} are keyless. \texttt{redemptions.py} and
\texttt{dai\_contagion.py} also require an Etherscan key. The key is read from the
environment only and is never written to any artifact. Re-running any of the six against
live data contacts an external API and overwrites the corresponding frozen file. The
reproduction sequence above does not require this step and should not begin with it.

The FIFO burn attribution underlying \texttt{data/redemptions\_by\_wallet.csv} is not
reproducible offline. We freeze and hash the attribution's \emph{output}, a 776-row
wallet-to-volume table. We do not freeze its input, the raw on-chain transfer logs the
attribution FIFO-matches. \texttt{redemptions.py --from-frozen} re-validates that frozen
output against its own manifest record. It does not repeat the attribution. The \$3.1B
PSM figure (Appendix~\ref{app:dai}) and the two public label sets
(Appendix~\ref{app:hhi}) are reporting-based in the same way: we freeze and hash what
they state, and the pipeline does not recompute them.

\begin{table}[t]
\centering
\caption{Claim map, first half. Each reported quantity, the results file it regenerates
from, and the key within that file. Both the file and the key are checked against the
results files themselves, and the section column is a cross-reference rather than a
typed number, so neither can drift.}
\label{tab:claimmap}
\scriptsize
\checkwidth{\scriptsize% Source: claim_map.yaml, via tables.py::write_tab_claimmap.
% Do not hand-edit: regenerate with `python tables.py`.
\begingroup
\footnotesize
\setlength{\tabcolsep}{2pt}
\Urlmuskip=0mu\relax
\def\UrlBreaks{\do\_\do\.\do\,\do\-}
\urlstyle{tt}
\begin{tabular}{@{}>{\raggedright\arraybackslash}p{3.4cm}c>{\raggedright\arraybackslash}p{2.9cm}>{\raggedright\arraybackslash}p{4.0cm}@{}}
\toprule
Claim & Sec. & Results file & Key \\
\midrule
Worst-case floor and implied unremediated-loss probability & \ref{sec:ident} & \path{number1b_rational_bound.json} & \path{p_obs}, \path{floor_worstcase}, \path{implied_failure_prob_rho0} \\
Non-fundamental share, lower edge and band & \ref{sec:ident} & \path{number1_nonfundamental_share.json} & \path{headline_band.low}, \path{headline_band.high} \\
Denominator sensitivity of phi and the share & \ref{sec:ident} & \path{phi_denominator_sensitivity.json} & \path{phi_conventions}, \path{grid}, \path{feed_trough_only_at_phi_0.08} \\
Dated impossibility window and per-coin firing counts & \ref{sec:ident} & \path{impossibility_window.json} & \path{impossibility_region_usdc}, \path{coins.USDC.independence.effective_n} \\
Dilution-adjusted phi(t) and its interpolation band & \ref{sec:sensitivity} & \path{phi_dynamic.json} & \path{variant_i}, \path{interpolation_band}, \path{firing_set} \\
Cross-feed trough corroboration and below-floor depth & \ref{sec:sensitivity} & \path{trough_corroboration.json} & \path{firing_report}, \path{below_floor_depth}, \path{independence} \\
Hourly observation counts per coin & \ref{app:repro} & \path{hourly_counts.json} & \path{per_coin}, \path{window.expected_hours} \\
Cross-coin placebo troughs and break threshold & \ref{sec:placebo} & \path{number3_placebo.json} & \path{USDC}, \path{DAI}, \path{USDT}, \path{break_threshold} \\
USDT premium depth and timing & \ref{sec:placebo} & \path{usdt_premium.json} & \path{max_premium_bps}, \path{premium_volume_usdt} \\
USDT premium peak timestamps by venue & \ref{sec:placebo} & \path{usdt_premium_timing.json} & \path{peaks.coinbase_usdt_usd}, \path{peaks.composite_usdt} \\
Redemption concentration and entity resolution & \ref{app:hhi} & \path{number2_redemption_hhi.json} & \path{upper_bound.hhi_raw}, \path{entity_resolved}, \path{coverage} \\
\bottomrule
\end{tabular}
\endgroup
}
% Source: claim_map.yaml, via tables.py::write_tab_claimmap.
% Do not hand-edit: regenerate with `python tables.py`.
\begingroup
\footnotesize
\setlength{\tabcolsep}{2pt}
\Urlmuskip=0mu\relax
\def\UrlBreaks{\do\_\do\.\do\,\do\-}
\urlstyle{tt}
\begin{tabular}{@{}>{\raggedright\arraybackslash}p{3.4cm}c>{\raggedright\arraybackslash}p{2.9cm}>{\raggedright\arraybackslash}p{4.0cm}@{}}
\toprule
Claim & Sec. & Results file & Key \\
\midrule
Worst-case floor and implied unremediated-loss probability & \ref{sec:ident} & \path{number1b_rational_bound.json} & \path{p_obs}, \path{floor_worstcase}, \path{implied_failure_prob_rho0} \\
Non-fundamental share, lower edge and band & \ref{sec:ident} & \path{number1_nonfundamental_share.json} & \path{headline_band.low}, \path{headline_band.high} \\
Denominator sensitivity of phi and the share & \ref{sec:ident} & \path{phi_denominator_sensitivity.json} & \path{phi_conventions}, \path{grid}, \path{feed_trough_only_at_phi_0.08} \\
Dated impossibility window and per-coin firing counts & \ref{sec:ident} & \path{impossibility_window.json} & \path{impossibility_region_usdc}, \path{coins.USDC.independence.effective_n} \\
Dilution-adjusted phi(t) and its interpolation band & \ref{sec:sensitivity} & \path{phi_dynamic.json} & \path{variant_i}, \path{interpolation_band}, \path{firing_set} \\
Cross-feed trough corroboration and below-floor depth & \ref{sec:sensitivity} & \path{trough_corroboration.json} & \path{firing_report}, \path{below_floor_depth}, \path{independence} \\
Hourly observation counts per coin & \ref{app:repro} & \path{hourly_counts.json} & \path{per_coin}, \path{window.expected_hours} \\
Cross-coin placebo troughs and break threshold & \ref{sec:placebo} & \path{number3_placebo.json} & \path{USDC}, \path{DAI}, \path{USDT}, \path{break_threshold} \\
USDT premium depth and timing & \ref{sec:placebo} & \path{usdt_premium.json} & \path{max_premium_bps}, \path{premium_volume_usdt} \\
USDT premium peak timestamps by venue & \ref{sec:placebo} & \path{usdt_premium_timing.json} & \path{peaks.coinbase_usdt_usd}, \path{peaks.composite_usdt} \\
Redemption concentration and entity resolution & \ref{app:hhi} & \path{number2_redemption_hhi.json} & \path{upper_bound.hhi_raw}, \path{entity_resolved}, \path{coverage} \\
\bottomrule
\end{tabular}
\endgroup

\end{table}

\begin{table}[t]
\centering
\caption{Claim map, second half. Columns as in Table~\ref{tab:claimmap}.}
\label{tab:claimmapb}
\scriptsize
\checkwidth{\scriptsize% Source: claim_map.yaml, via tables.py::write_tab_claimmap.
% Do not hand-edit: regenerate with `python tables.py`.
\begingroup
\footnotesize
\setlength{\tabcolsep}{2pt}
\Urlmuskip=0mu\relax
\def\UrlBreaks{\do\_\do\.\do\,\do\-}
\urlstyle{tt}
\begin{tabular}{@{}>{\raggedright\arraybackslash}p{3.4cm}c>{\raggedright\arraybackslash}p{2.9cm}>{\raggedright\arraybackslash}p{4.0cm}@{}}
\toprule
Claim & Sec. & Results file & Key \\
\midrule
FIFO burn-attribution rule & \ref{app:hhi} & \path{fifo_rule.json} & \path{matching_window}, \path{tie_break}, \path{truncation} \\
DAI reported USDC exposure & \ref{app:dai} & \path{number4_dai_contagion.json} & \path{psm_share_of_dai_supply}, \path{source_collateral_share}, \path{no_decomposition} \\
Specificity panel rows and breach combinations & \ref{sec:sensitivity} & \path{specificity_panel.json} & \path{row1_usdc_march_2023}, \path{row2_tether_2019}, \path{row4_terra_2022} \\
Terra degenerate control & \ref{sec:sensitivity} & \path{second_event_terra.json} & \path{phi_terra_exposure}, \path{degenerate_caveat} \\
Excluded-coin calm-baseline evidence & \ref{app:excl} & \path{placebo_excluded_coins.json} & \path{USDP}, \path{GUSD}, \path{BUSD} \\
Exposure-ordering null probability & \ref{sec:placebo} & \path{ordering_null.json} & \path{probability_as_fraction} \\
TrueUSD scope-condition deviation & \ref{sec:disc} & \path{tusd_2023.json} & \path{deviation_from_par_bps} \\
Make-whole pledge timing bound & \ref{sec:ident} & \path{makewhole_timing.json} & \path{circle_makewhole_pledge.by_trough_reference} \\
Time-value discount on a delayed redemption & \ref{sec:disc} & \path{time_value.json} & \path{discount_bps_range_2to3d} \\
Trough print volume & \ref{sec:placebo} & \path{trough_volume.json} & \path{volume_in_trough_hour_usdc} \\
Hour-by-hour panel series for 11 March & \ref{sec:ident} & \path{panel_11mar.json} & \path{panel}, \path{cross_coin_cross_window_summary} \\
\bottomrule
\end{tabular}
\endgroup
}
% Source: claim_map.yaml, via tables.py::write_tab_claimmap.
% Do not hand-edit: regenerate with `python tables.py`.
\begingroup
\footnotesize
\setlength{\tabcolsep}{2pt}
\Urlmuskip=0mu\relax
\def\UrlBreaks{\do\_\do\.\do\,\do\-}
\urlstyle{tt}
\begin{tabular}{@{}>{\raggedright\arraybackslash}p{3.4cm}c>{\raggedright\arraybackslash}p{2.9cm}>{\raggedright\arraybackslash}p{4.0cm}@{}}
\toprule
Claim & Sec. & Results file & Key \\
\midrule
FIFO burn-attribution rule & \ref{app:hhi} & \path{fifo_rule.json} & \path{matching_window}, \path{tie_break}, \path{truncation} \\
DAI reported USDC exposure & \ref{app:dai} & \path{number4_dai_contagion.json} & \path{psm_share_of_dai_supply}, \path{source_collateral_share}, \path{no_decomposition} \\
Specificity panel rows and breach combinations & \ref{sec:sensitivity} & \path{specificity_panel.json} & \path{row1_usdc_march_2023}, \path{row2_tether_2019}, \path{row4_terra_2022} \\
Terra degenerate control & \ref{sec:sensitivity} & \path{second_event_terra.json} & \path{phi_terra_exposure}, \path{degenerate_caveat} \\
Excluded-coin calm-baseline evidence & \ref{app:excl} & \path{placebo_excluded_coins.json} & \path{USDP}, \path{GUSD}, \path{BUSD} \\
Exposure-ordering null probability & \ref{sec:placebo} & \path{ordering_null.json} & \path{probability_as_fraction} \\
TrueUSD scope-condition deviation & \ref{sec:disc} & \path{tusd_2023.json} & \path{deviation_from_par_bps} \\
Make-whole pledge timing bound & \ref{sec:ident} & \path{makewhole_timing.json} & \path{circle_makewhole_pledge.by_trough_reference} \\
Time-value discount on a delayed redemption & \ref{sec:disc} & \path{time_value.json} & \path{discount_bps_range_2to3d} \\
Trough print volume & \ref{sec:placebo} & \path{trough_volume.json} & \path{volume_in_trough_hour_usdc} \\
Hour-by-hour panel series for 11 March & \ref{sec:ident} & \path{panel_11mar.json} & \path{panel}, \path{cross_coin_cross_window_summary} \\
\bottomrule
\end{tabular}
\endgroup

\end{table}

An anonymized artifact repository is available at
\url{https://anonymous.4open.science/r/usdc-depeg-artifact-F62D/}.

